\chapter{State of the Art and Research Gap}
\label{ch:sota}

\section{LLM-based evaluation as a measurement problem}

LLM-based evaluation has become an important practical response to the cost of expert annotation. A language model can apply a rubric, compare responses, and provide structured output. The literature, however, does not present this as a solved replacement for human evaluation. It presents a family of methods whose quality depends on task, prompt, model, reference material, and protocol. The appropriate question is therefore not whether an LLM can judge in the abstract, but which reliability properties are observed under a reproducible design.

The MT-Bench and Chatbot Arena study by \citet{zheng2023judging} made this issue concrete by using strong models as evaluators while examining agreement, position effects, verbosity effects, and self-enhancement concerns. It also established MT-Bench as a useful multi-turn setting for comparative evaluation. The present project adopts the need for a controlled and traceable benchmark context, while focusing on a smaller, locally reproducible set of stored pairwise decisions and explicitly separating 80 distinct questions from their 160 turn-level prompt records.

\section{Alignment with human judgements}

Several approaches seek stronger correspondence with human assessment. G-EVAL uses a structured prompting procedure with chain-of-thought and form-filling ideas to evaluate natural-language generation, reporting improved correlation with human judgements in its target settings \citep{liu2023geval}. This direction motivates explicit rubric and output-parsing design. Yet correlation or agreement alone does not establish that an evaluator is robust to order, repetition, or presentation changes. A judge can be aligned on average but still produce unstable or protocol-sensitive individual decisions.

For that reason, the present project treats human labels as \emph{human preference reference labels}. They are the defined comparison target for RQ1, not an assertion of objective correctness. The three-class outcome space includes answer 1, answer 2, and tie. This is an important distinction: collapsing ties or treating all disagreement as one homogeneous type can make a simple percentage look more decisive than the evaluation setting warrants.

\section{Position and presentation sensitivity}

Order sensitivity is among the most visible risks in pairwise evaluation. \citet{shi2025judging} systematically study position bias across judges, tasks, and candidate models, distinguishing repetition stability, position consistency, and preference fairness. \citet{wang2024fair} likewise report fairness concerns for large-language-model evaluators and study a calibration-oriented response. These studies motivate a design that does not infer order sensitivity from a single displayed order. Instead, the project evaluates AB and BA passes and maps both decisions to the original answer identities before calculating paired outcome disagreement or decisive flips.

Presentation sensitivity can be confused with content sensitivity. The literature on length control provides a related example: \citet{dubois2024length} argue that automatic evaluators can be confounded by answer length and propose a controlled adjustment procedure. Their work motivates experimental control, but it does not license a universal inference from one perturbation. RQ4 therefore asks only whether a specifically frozen redundant-text variant receives a stable win in the present population. RQ5 similarly asks only whether a frozen list-prefix/presentation treatment receives a stable win. Neither estimand claims to settle all effects of answer length or formatting.

\section{Mitigation and aggregation}

Mitigation approaches vary in what they change. Some modify prompts or scoring rules; others filter observations; others aggregate multiple judges. The present project distinguishes two families because their outputs are not directly comparable as if they had the same units and comparator. DUAL\_SWAP is a within-judge presentation-consistency filter. It retains cases satisfying a fixed mapped-order criterion and therefore changes coverage. Strict multi-judge consensus aggregates independent judge votes and compares consensus against an equal-weight individual-judge baseline on the same retained pairs. It is a secondary RQ7 strategy, not a replacement for DUAL\_SWAP.

This distinction matters scientifically. A multi-judge agreement change measured on a strict four-judge retained population cannot be subtracted from a DUAL\_SWAP agreement change measured on a different retained matched population to claim superiority. Both are valuable results, but they answer different operational questions. The report therefore preserves the comparison boundary in its tables and interface description.

\section{Research gap}

The literature identifies the central risks, but a research project also needs to show how those risks are managed end to end. The gap addressed here is an integrated, reproducible workflow that joins: (i) exact source-text provenance, (ii) controlled RQ-specific designs, (iii) stored operational terminal states, (iv) analysis-run-pinned API values that fail closed when a required artifact is absent or ambiguous, and (v) cautious claims tied to a named frozen estimator. The project also makes a practical contribution by separating exploratory historical telemetry, final controlled evidence, and live manual trials so that a convenient interface cannot silently contaminate a final conclusion.

\section{Conceptual framework}

\begin{figure}[H]
\centering
\begin{tikzpicture}[node distance=8mm and 9mm,>=Latex, every node/.style={font=\small}, box/.style={draw,rounded corners,align=center,minimum height=10mm,minimum width=30mm,fill=gray!7}]
\node[box] (data) {Canonical source text\\and human reference labels};
\node[box,right=of data] (protocol) {Frozen controlled\\protocols};
\node[box,right=of protocol] (obs) {Persisted judge\\observations};
\node[box,below=of protocol] (analysis) {Pinned analysis runs\\and uncertainty};
\node[box,right=of analysis] (claims) {Cautious scientific\\claims and UI};
\draw[->] (data)--(protocol); \draw[->] (protocol)--(obs); \draw[->] (obs)--(analysis); \draw[->] (analysis)--(claims); \draw[->] (data)--(analysis);
\end{tikzpicture}
\caption{The project treats reliability claims as a chain from exact source text through a frozen protocol and persisted observations to a pinned analysis.}
\label{fig:conceptual-framework}
\end{figure}

The framework is intentionally conservative. Every link can fail independently: source text may be mismatched, a provider response may be invalid, an observation may be terminally unavailable, a repeated pass may be incomplete, or an analysis identity may be missing. A reliable system records these conditions and fails closed rather than silently manufacturing a result. This design guides the rest of the report.
